\section{Introduction}

Entropy is a fundamental concept that is directly correlated with the security of the cryptographic system.
A key with high entropy indicates a high level of difficulty for an attacker to recover the key from the ciphertext.
However, in practice, modern cryptographic protocols, such as WPA2-PSK, Bluetooth Low Energy~(BLE), and Bluetooth Classic~(BC), usually tolerate low-entropy keys to accommodate user-friendly designs, maintain backward compatibility, or support resource-constrained devices.
This raises a security issue concerning the recovery of low-entropy keys in these protocols, as demonstrated by attacks like PMKID~\cite{pmkidattack} and KNOB~\cite{antonioli2019knob}.

The use of symbolic models to formally analyze protocol vulnerabilities, pioneered by Dolev et al.~\cite{dolev1983security}, has become a mainstream research approach.
Numerous tools have been developed to automatically analyze protocols and identify potential attacks, including AVISPA\cite{armando2005avispa}, ProVerif~\cite{blanchet2018proverif}, Maude-NPA~\cite{escobar2007maude}, Scyther~\cite{cremers2008scyther}, \Tamarin{}~\cite{meier2013tamarin}, and others.
In symbolic models, messages and keys are abstracted as terms in term algebra, rather than as concrete bit sequences as used in computational models.
Consequently, the entropy of keys is lost during this abstraction process, making it impossible for existing symbolic model tools to distinguish low-entropy keys or analyze the brute-force capabilities of attackers.

This raises a critical question: How to capture the brute-force behaviors of attackers within symbolic models?
Some studies\cite{shi2023blesc, claverie2023tamarin} have proposed solutions by manually identifying keys that are vulnerable to compromise and defining specific rules that allow these keys to be directly revealed to attackers.
However, these studies do not provide a general or automated method for analyzing the impact of low-entropy keys on the entire protocol.
Firstly, this method relies on the manual identification of key vulnerabilities, limiting the scalability and objectivity of the analysis.
Secondly, manual analysis is inherently constrained, as it only offers attacks a limited set of predefined ways to compromise low-entropy keys.
As a result, it fails to provide a comprehensive strategy for effectively exploiting such vulnerabilities.

To automatically analyze protocols tolerating low-entropy keys, we encounter three main challenges:
\begin{enumerate}[leftmargin=*]
    \item \textbf{Entropy Representation:}
    Symbolic analysis tools typically preserve only the logical structure of messages or keys, making it impossible to deduce and analyze the entropy of keys.
    Additional mechanisms are needed to describe a key's entropy in order to cover general brute-force attacks within the symbolic model.
    \item \textbf{Prerequisites for Brute-Force Attacks:}
    Brute-force attacks involve an attacker repeatedly attempting keys, relying on prior knowledge of the algorithm, its outputs, and other relevant inputs beside targeted keys~(e.g., salt values, initialization vectors). 
    Accurately and generally capturing these prerequisite conditions within a symbolic model presents a significant challenge, as these conditions vary across different protocols and keys.
    \item \textbf{Enhanced Adversarial Capabilities:}
    In existing works, an attacker is typically limited to compromising specific keys. 
    However, an attacker can attempt to brute-force any key as long as its entropy doesn't exceed its computational capacity. 
    The attacker can also deconstruct complex keys into simpler sub-keys and brute-force them individually. 
    A challenge lies in modeling this flexible brute-forcing capability within the symbolic model.
\end{enumerate}

This paper aims to bridge the existing gap by introducing a general and automated method for modeling low-entropy keys and the adversarial brute-force capabilities within symbolic analysis frameworks. 
We focus on enhancing \Tamarin{} to describe general brute-force attacks conducted by adversaries. 
\Tamarin{} is a state-of-the-art protocol verification tool widely used for the formal analysis of many protocols, including TLS\cite{cremers2017tls13}, 5G-AKA\cite{basin20185G, wang20215G}, EMV\cite{basin2021emv,radu2022emv}, etc.
We have designed a new syntax, which is a superset of the syntax of \Tamarin{}, allowing for the labeling of low-entropy variables in symbolic models. 
Additionally, we have developed an automated protocol analysis tool supporting such low-entropy annotated models, which we call EntropyVerif\cite{toolrepo}. 
This tool supports an adversary model that extends beyond the Dolev-Yao model, capable of capturing all brute-force behaviors against security protocols, including complex attacks where an attacker synthesizes and utilizes various attack capabilities.
Our approach can be used to identify security issues arising from low-entropy keys in protocols and to verify the security of protocols that tolerate low-entropy keys.

\smallskip
\noindent\textbf{A new attack on Bluetooth Classic.}
Applying EntropyVerif to BC session establishment immediately exposes a new protocol-level attack that we call the Unilateral Multi-Downgrade (UMD) attack.
The BC specification allows the two endpoints to negotiate a session-key length between 1 and 16 bytes.
UMD exploits the fact that when one side is repeatedly renegotiated down to a short key, the resulting session keys are not independent fresh randomness but correlated fragments derived from a shared long-term secret; an attacker can therefore brute-force each fragment separately and stitch the full-length session key back together.
To our knowledge this is the first protocol-level attack on BC discovered by a symbolic verifier; Section~\ref{sec:umd-attack} gives the full complexity analysis and the attack trace.
The attack has been responsibly disclosed to the Bluetooth SIG.

\noindent\textbf{Contributions.}
Our primary contributions are as follows:
\begin{itemize}[leftmargin=*]
    \item \emph{Protocol-agnostic methodology.} We relax the \emph{perfect encryption assumption} in the symbolic model and introduce the \texttt{MDerive} message-derivation fact together with a three-tier breaking-oracle design (value-breaking, relationship-breaking, message-derivation). Unlike prior ad-hoc extensions~\cite{shi2023blesc,claverie2023tamarin,jangid2023bluepe} that hard-code reveal rules for a specific protocol, our abstractions are function-symbol agnostic and support recursive decomposition of compound keys.
    \item \emph{Tool without kernel modification.} We developed EntropyVerif, a syntax-sugar extension of \Tamarin{}, in which users annotate low-entropy terms with \texttt{:= L}.
          \emph{The compiler emits multiset-rewriting rules so that verification soundness is inherited directly from the unmodified \Tamarin{} backend.}
    \item \emph{Rediscovery of known attacks.} Using EntropyVerif we automatically reproduce the PMKID attack on WPA2-PSK, the KNOB attack and the Method Confusion attack on BLE pairing, without protocol-specific reveal rules.
    \item \emph{A new attack on Bluetooth Classic (UMD).} Beyond rediscovery, the compound-key decomposition naturally supported by \texttt{MDerive} leads the verifier to a previously unknown protocol-level attack: the Unilateral Multi-Downgrade attack on BC, which reduces 128-bit session-key recovery from $2^{128}$ to $n \cdot 2^{128/n}$, i.e.\ as low as $2^{12}$ when $n = 16$---a factor of $2^{116}$ reduction. The attack has been responsibly disclosed to the Bluetooth SIG.
\end{itemize}

\noindent\textbf{Overview.} 
% Section \ref{sec:2background} reviews protocols that tolerate low-entropy keys and the protocol analysis tool \Tamarin{}.
Section \ref{sec:2background} reviews the definition of entropy and brute-force attacks, protocols that tolerate low-entropy keys, and the protocol analysis tool \Tamarin{}.
Section \ref{sec:3symbolic-modeling} introduces our modeling approach for low-entropy analysis, detailing the key-breaking oracles. 
Section \ref{sec:4extension-and-compilation} discusses EntropyVerif, our extension of \Tamarin{}.
Section \ref{sec:5case-study} presents case studies where we apply our method to analyze real-world protocols. 
Finally, in Section \ref{sec:6results}, we present the results of our case studies, where we discover known vulnerabilities and a novel attack. 
We offer a review of related work in Section \ref{sec:7related-work} and provide a conclusion and discussion of our work in Section \ref{sec:8discussion-conclusion}.
